\section*{Chapitre \ref{ch:iv1:regression} --- Régression}
\markboth{\protect\omTocMark{}{Réponses~: Régression}}{\protect\omTocMark{}{Réponses~: Régression}}

\begin{solution}{iq:iv1:regression:1}
La pente inverse peut prendre toute valeur entre 0 (exclu) et $1/2$. Une pente inverse de $0.18$ correspond à une corrélation de $0.6$.

Avec $b_1 = \mathrm{Cov}(x,y)/\mathrm{Var}(x)$ et $b_2 = \mathrm{Cov}(x,y)/\mathrm{Var}(y)$, le produit vaut
\[
b_1 b_2 = \frac{\mathrm{Cov}(x,y)^2}{\mathrm{Var}(x)\mathrm{Var}(y)} = \rho^2 .
\]
Donc $b_2 = \rho^2/2$, avec $\rho^2$ dans $(0, 1]$. Les deux pentes ont toujours le même signe. Si $b_2 = 0.18$, alors $\rho^2 = 0.36$, et la corrélation vaut $0.6$, positive puisque les pentes le sont.

\iqwrong{«~$b_2 = 1/b_1 = 1/2$.~» C'est vrai seulement avec $|\rho| = 1$. Chaque régression minimise les erreurs verticales selon son propre axe, donc les deux droites diffèrent, et elles ne coïncident que si tous les points sont alignés.}
\iqwhy{Le sens de la régression n'est pas indifférent lorsqu'on fixe un ratio de couverture. Les deux régressions donnent des ratios différents, et le bon dépend de ce que l'on veut minimiser.}
\end{solution}

\begin{solution}{iq:iv1:regression:2}
Les coefficients et le $R^2$ ne changent pas. Les erreurs types diminuent d'un facteur d'environ $\sqrt2$, et les statistiques $t$ augmentent d'environ $\sqrt2$.

$X^\top X$ et $X^\top y$ doublent tous deux, donc $\hat\beta = (X^\top X)^{-1}X^\top y$ est inchangé. Les résidus se répètent, donc la RSS double tandis que le nombre de degrés de liberté double presque, et $s^2$ reste à peu près le même. Comme $(X^\top X)^{-1}$ est divisé par deux, l'estimation de la variance est divisée par deux. Exactement, avec $n$ lignes et $p$ coefficients, les erreurs types évoluent d'un facteur $\sqrt{(n-p)/(2n-p)}$, ce qui pour $n = 250$ et $p = 2$ donne $0.706$, à comparer avec $1/\sqrt2 = 0.707$.

\iqwhy{Rien n'a été appris, et pourtant la statistique $t$ a augmenté. Des données qui se chevauchent produisent le même effet~: des rendements sur 20 jours échantillonnés chaque jour ressemblent à 20 fois plus de données, alors qu'ils ne contiennent qu'environ une observation indépendante tous les 20 jours. Avec un signal lent, la statistique $t$ est surestimée d'environ $\sqrt{20} \approx 4.5$. Il faut utiliser des échantillons sans recouvrement ou des erreurs types qui tiennent compte du chevauchement.}
\end{solution}

\begin{solution}{iq:iv1:regression:3}
Vendre à découvert $1.6$ dollars d'indice par dollar d'action. La position couverte a une volatilité de 18~\%. En utilisant la régression inverse, on vendrait $2.5$ dollars et la volatilité qui resterait serait de $22.5~\%$.

\begin{sloppypar}La variance de l'action moins $h$ fois l'indice est un trinôme en $h$, minimisée en $h = \mathrm{Cov}/\mathrm{Var}(\text{indice}) = \rho\,\sigma_s/\sigma_i = 0.8 \cdot 30/15 = 1.6$. C'est la pente de la régression de l'action sur l'indice. La volatilité résiduelle vaut $\sigma_s\sqrt{1 - \rho^2} = 30 \cdot 0.6 = 18~\%$, de sorte que la couverture élimine $\rho^2 = 64~\%$ de la variance.\end{sloppypar}

La régression de l'indice sur l'action a pour pente $\rho\,\sigma_i/\sigma_s = 0.4$, et son inverse donne $2.5$. La variance vaut alors $30^2 + 2.5^2 \cdot 15^2 - 2 \cdot 2.5 \cdot 360 = 506.25$, soit une volatilité de $22.5~\%$. La couverture est alors excessive, d'un facteur $1/\rho^2$.

\iqwhy{C'est une décision quotidienne sur un desk. Notons aussi qu'un ratio estimé sur des rendements bruités a tendance à sortir faible, et non élevé, car le bruit dans la variable explicative ramène la pente vers zéro.}
\end{solution}

\begin{solution}{iq:iv1:regression:4}
Avec une constante, $R^2$ appartient à $[0, 1]$, et $R^2 = 1$ signifie que tous les résidus sont nuls. Ajouter une variable explicative ne peut jamais le faire baisser. Sans constante, cette formule peut donner une valeur très inférieure à zéro.

Avec une constante, la somme des carrés totaux se décompose en somme expliquée plus somme des résidus, car les résidus sont orthogonaux aux valeurs ajustées et de somme nulle\index{equations normales@équations normales}. Ainsi $R^2$ est un rapport compris entre 0 et 1, et il est égal au carré de la corrélation entre $y$ et les valeurs ajustées. Une nouvelle variable explicative peut reproduire l'ajustement précédent en prenant un coefficient nul, donc la RSS ne peut pas augmenter.

Sans constante, l'ajustement n'est pas tenu de faire aussi bien que la moyenne. Si $y$ a une moyenne élevée et que $x$ est centrée sur zéro, la RSS est proche de $\sum y_i^2$, bien supérieure à $\sum(y_i - \bar y)^2$, et la formule donne un grand nombre négatif.

\iqwrong{Lire un $R^2$ plus élevé après l'ajout d'une variable explicative comme une preuve. Il augmente même pour du pur bruit.}
\iqfollow{Et hors échantillon (out-of-sample)~?}{Il peut être négatif pour n'importe quel modèle, lorsque celui-ci prédit moins bien que la moyenne.}
\end{solution}

\begin{solution}{iq:iv1:regression:5}
La statistique $t$ vaut $0.4$, et il faudrait 100 ans pour atteindre une statistique $t$ de 2.

L'erreur type de la constante vaut environ la volatilité résiduelle divisée par la racine carrée du nombre d'années, soit $10\%/\sqrt4 = 5\%$ par an. Donc $t = 2/5 = 0.4$. En termes de ratio d'information\index{ratio d'information}, c'est-à-dire l'alpha divisé par la volatilité résiduelle, ici $0.2$, on a $t = \text{ratio} \cdot \sqrt{\text{années}}$. Pour atteindre 2, il faut $\sqrt{\text{années}} = 2/0.2 = 10$.

\iqfollow{Et si le ratio d'information était de 1~?}{Alors $t = 2$ est atteint en 4 ans.}
\iqwhy{Un alpha de cette taille ne peut pas être confirmé par un historique de performance sur une vie professionnelle. En pratique, on s'appuie sur la logique économique de la stratégie, sur une fréquence bien plus élevée de paris indépendants, ou sur des preuves en coupe transversale. La performance seule ne suffit pas.}
\end{solution}

\enlargethispage{-6\baselineskip}%
\begin{solution}{iq:iv1:regression:6}
La statistique $t$ vaut environ $3.54$. Le signal peut être précieux, avec un ratio de Sharpe annuel d'environ $1.1$.

Une régression simple vérifie l'identité exacte
\[
t^2 = \frac{R^2\,(n - 2)}{1 - R^2},
\]
donc $t = \sqrt{0.005 \cdot 2498/0.995} \approx 3.54$. Un $R^2$ faible associé à un $t$ élevé est normal~: $t$ croît avec la taille de l'échantillon, alors que $R^2$ ne le fait pas.

La corrélation avec le rendement vaut $\sqrt{0.005} \approx 0.071$. Détenir une position proportionnelle au signal standardisé $x$ rapporte $xy$, de moyenne environ $\rho\,\sigma_y$ et d'écart type environ $\sigma_y$, donc le ratio de Sharpe quotidien vaut environ $\rho$. En annualisant, on obtient $0.071 \cdot \sqrt{252} \approx 1.12$. La même réponse résulte de $t/\sqrt{\text{années}} = 3.54/\sqrt{9.9}$.

\iqwrong{«~Un $R^2$ d'un demi~pour~cent signifie qu'il n'explique rien.~» Sur des rendements quotidiens, c'est un signal fort, avant coûts.}
\iqwhy{Il faut juger la prédiction des rendements au ratio de Sharpe, après coûts. Ce calcul suppose aussi que les données ne sont pas issues d'une recherche parmi de nombreux signaux, et c'est en général le premier point à vérifier.}
\end{solution}

\begin{solution}{iq:iv1:regression:7}
L'estimateur est $\hat\beta = (X^\top X)^{-1}X^\top y$, et il existe de façon unique lorsque $X$ est de rang plein en colonnes.

On minimise $L(\beta) = \lVert y - X\beta\rVert^2 = y^\top y - 2\beta^\top X^\top y + \beta^\top X^\top X\beta$. En annulant le gradient $-2X^\top y + 2X^\top X\beta$, on obtient les équations normales\index{equations normales@équations normales} $X^\top X\beta = X^\top y$. La hessienne $2X^\top X$ est semi-définie positive, donc il s'agit d'un minimum, et celui-ci est unique si $X^\top X$ est inversible, ce qui revient au rang plein en colonnes. Avec moins de lignes que de colonnes, ou avec une colonne répétée, il existe de nombreuses solutions.

Géométriquement, $X\hat\beta$ est la projection orthogonale de $y$ sur l'espace engendré par les colonnes de $X$. Les résidus sont donc orthogonaux à chaque colonne, et, avec une constante, ils ont une somme nulle.

Ne pas inverser. Le conditionnement de $X^\top X$ est le carré de celui de $X$, si bien que l'inversion fait perdre deux fois plus de chiffres. Il faut résoudre par QR sur $X$, ou par Cholesky sur $X^\top X$ si $X$ est bien conditionnée.

Si $X$ ne tient pas en mémoire, on le lit par blocs et l'on accumule la matrice $p \times p$ $X^\top X$ et le vecteur $X^\top y$. Ce sont des sommes sur les lignes, donc on obtient la réponse exacte avec une mémoire proportionnelle à $p^2$.

\iqfollow{Et lorsque les lignes arrivent une à une~?}{Il faut mettre à jour l'inverse par moindres carrés récursifs, et ajouter un facteur d'oubli si la relation dérive.}
\end{solution}

\begin{solution}{iq:iv1:regression:8}
La procédure B donne exactement le coefficient de la régression complète. La procédure A donne ici $0.36$, et non 1.

C'est le théorème de Frisch--Waugh--Lovell\index{théorème de Frisch--Waugh--Lovell}~: le coefficient de $x$ dans la régression complète est égal à la pente des résidus de $y$ sur les résidus de $x$, après avoir retiré $z$ des deux. Retirer $z$ de $y$ seulement laisse $x$ porter encore sa part en $z$, qui n'est pas liée aux résidus de $y$ et dilue la pente. En notant $\tilde x$ la part de $x$ non expliquée par $z$,
\[
\text{pente de A} = \beta \cdot \frac{\mathrm{Var}(\tilde x)}{\mathrm{Var}(x)} = \beta\,(1 - \rho^2) = 1 \cdot (1 - 0.64) = 0.36 .
\]
\iqwrong{«~Retirer le marché des rendements, puis tester le signal sur les résidus.~» C'est une démarche courante, et c'est la procédure A. Elle biaise le coefficient du signal vers zéro d'un facteur $1 - \rho^2$, et plus fortement si le signal est fortement corrélé au facteur.}
\iqwhy{Il faut neutraliser le signal autant que les rendements. Si l'on neutralise seulement les rendements par secteur et que le signal présente un biais sectoriel, on sous-estime ce que le signal rapporterait après couverture.}
\end{solution}

\begin{solution}{iq:iv1:regression:9}
La pente vaut $-0.5$. Le signe est faux~: le vrai coefficient de $x$ vaut $+1$.

Il s'agit d'un biais de variable omise\index{biais de variable omise}. La pente de la régression courte vaut $\mathrm{Cov}(x, y)/\mathrm{Var}(x)$. En remplaçant le vrai modèle,
\[
\frac{\mathrm{Cov}(x,\, x - 3z + \varepsilon)}{\mathrm{Var}(x)} = 1 - 3 \cdot \frac{\mathrm{Cov}(x, z)}{\mathrm{Var}(x)} = 1 - 3 \cdot 0.5 = -0.5 .
\]
En général, la pente courte est le vrai coefficient, auquel s'ajoute le coefficient de $z$ multiplié par la pente de $z$ sur $x$. Le biais ne disparaît que si $z$ n'a aucun effet, ou si $x$ et $z$ ne sont pas corrélés, et il ne diminue pas avec davantage de données.

\iqfollow{Et si $x$ et $z$ ne sont pas corrélés~?}{La pente vaut 1, sans biais. La variable $z$ omise ne fait qu'ajouter du bruit.}
\iqwhy{Un signal corrélé au bêta de marché ou au secteur capte leurs rendements. Avant de croire à un signal, il faut le régresser sur les facteurs connus et examiner ce qui reste.}
\end{solution}

\begin{solution}{iq:iv1:regression:10}
Du bruit dans $x$ donne une pente de $1/2$. Du bruit dans $y$ seul laisse la pente à 1. Avec du bruit dans les deux, la vraie pente se situe entre la pente de $y$ sur $x$ et l'inverse de la pente de $x$ sur $y$.

Il s'agit des erreurs sur les variables\index{erreurs sur les variables}. Le bruit sur le régresseur augmente $\mathrm{Var}(x)$ et laisse la covariance inchangée~:
\[
\text{pente} = \frac{\mathrm{Cov}(x^* + n,\, y)}{\mathrm{Var}(x^*) + \mathrm{Var}(n)} = \frac{1}{1 + 1} = \frac12 .
\]
Le facteur $\mathrm{Var}(x^*)/\mathrm{Var}(x)$ est l'atténuation. Le bruit sur $y$ passe dans le terme d'erreur~: il élargit les erreurs types mais n'introduit aucun biais.

Avec du bruit sur les deux variables, la pente de $y$ sur $x$ est trop plate et l'inverse de la pente de $x$ sur $y$ est trop raide. Avec les nombres ci-dessus et un bruit unitaire aussi sur $y$, ces valeurs valent $0.5$ et $2$, de part et d'autre de la vraie valeur 1.

\iqfollow{Existe-t-il une droite unique qui traite $x$ et $y$ de la même manière~?}{Les moindres carrés totaux\index{moindres carrés totaux} minimisent la distance perpendiculaire et retiennent la première direction principale. Leur pente est
$\bigl[(s_{yy} - s_{xx}) + \sqrt{(s_{yy} - s_{xx})^2 + 4s_{xy}^2}\,\bigr]/(2s_{xy})$, ici exactement égale à 1. La méthode suppose un bruit égal sur les deux variables et dépend de leurs unités, donc on standardise d'abord ou on utilise le rapport de bruit connu.}
\iqwhy{Les bêtas, volatilités et signaux estimés sont tous bruités, donc les régressions sur ces variables ressortent plates. Si l'on dimensionne une couverture à partir d'un régresseur bruité, on sous-couvre.}
\end{solution}

\begin{solution}{iq:iv1:regression:11}
Chaque erreur type est plus grande d'un facteur d'environ $2.29$. La somme des coefficients reste bien estimée, et oui, les deux statistiques $t$ peuvent être petites alors que la statistique $F$ est grande.

Le facteur d'inflation de la variance\index{facteur d'inflation de la variance} vaut $1/(1 - \rho^2) = 1/0.19 = 5.26$, et sa racine carrée vaut $2.29$.

Pour des régresseurs de variance unitaire, la covariance des coefficients est proportionnelle à l'inverse de $\bigl(\begin{smallmatrix}1 & \rho\\ \rho & 1\end{smallmatrix}\bigr)$. Cela donne $\mathrm{Var}(b_1 + b_2) \propto 2/(1 + \rho) = 1.05$ et $\mathrm{Var}(b_1 - b_2) \propto 2/(1 - \rho) = 20$. Les données fixent la somme 19 fois plus précisément que la différence, car les deux régresseurs bougent presque toujours ensemble. La difficulté réside dans la répartition de l'effet entre eux, ce qui se traduit par des coefficients qui changent de signe d'un échantillon à l'autre.

Le test $F$ porte sur l'effet conjoint, donc il n'en souffre pas. Avec 100 observations et $R^2 = 0.2$, la statistique de test vaut $F = (0.2/2)/(0.8/97) \approx 12.1$, bien au-dessus de la valeur critique à 5~\% d'environ $3.1$. Dans une simulation avec une corrélation $0.95$ et des coefficients vrais de $0.2$ chacun, plus de la moitié des échantillons ont à la fois $|t| < 2$, et la plupart de ceux-ci gardent quand même un $F$ significatif.

\iqwhy{Retirer l'une d'elles, les combiner ou utiliser une régression ridge. Si l'on n'a besoin que de prédictions et que la corrélation est stable, on peut ne rien faire, car la colinéarité ne nuit pas à l'ajustement.}
\end{solution}

\begin{solution}{iq:iv1:regression:12}
Les deux droites sont $y = x/2$ et $x = (1 + y)/2$. La corrélation vaut $1/2$.

Comme la région donne des moyennes conditionnelles simples, on peut écrire directement la régression. Sachant $x$, la hauteur $y$ est uniforme sur $[0, x]$, donc $E[y \mid x] = x/2$. Sachant $y$, la valeur $x$ est uniforme sur $[y, 1]$, donc $E[x \mid y] = (1 + y)/2$. Ce sont deux droites, et chacune est donc aussi la droite de régression théorique\index{espérance conditionnelle}, comme le montre la \cref{fig:iv1:triangle}. Les deux pentes valent $1/2$, et le produit des pentes est égal à $\rho^2$, donc $\rho = 1/2$.

\begin{omfigure}
\begin{tikzpicture}[x=4.2cm,y=4.2cm]
\fill[omDef!10] (0,0) -- (1,0) -- (1,1) -- cycle;
\draw[thick,omDef] (0,0) -- (1,0) -- (1,1) -- cycle;
\draw[->] (0,0) -- (1.2,0) node[right,font=\footnotesize] {$x$};
\draw[->] (0,0) -- (0,1.15) node[above,font=\footnotesize] {$y$};
\draw[very thick,omThm] (0,0) -- (1,0.5);
\draw[very thick,omExo] (0.5,0) -- (1,1);
\fill (2/3,1/3) circle (0.012);
\node[font=\footnotesize,anchor=north west] at (1.02,0.5) {$y = x/2$};
\node[font=\footnotesize,anchor=west] at (1.02,1) {$x = (1+y)/2$};
\node[font=\footnotesize,anchor=north west] at (0.69,0.32) {moyenne};
\node[below,font=\footnotesize] at (1,0) {$1$};
\node[left,font=\footnotesize] at (0,1) {$1$};
\end{tikzpicture}

\omcaption{Le triangle et ses deux droites de régression. Elles se coupent au point moyen $(2/3, 1/3)$. Tracées dans les mêmes axes, l'une a pour pente $1/2$ et l'autre la pente 2, de part et d'autre de la pente 1 de l'arête diagonale.}\label{fig:iv1:triangle}
\end{omfigure}

Pour vérifier par intégration, $E[x] = 2/3$, $E[y] = 1/3$, les deux variances valent $1/18$ et la covariance vaut $1/36$. Chaque pente vaut donc $(1/36)/(1/18) = 1/2$.

\iqwrong{Prendre l'arête diagonale, $y = x$, pour droite des données est une erreur. C'est la frontière de la région, et non une droite de régression.}
\end{solution}

\begin{solution}{iq:iv1:regression:13}
Les moindres carrés ordinaires n'ont pas de solution unique, car seule la somme des deux coefficients est identifiée. La régression ridge répartit l'effet à parts égales. Le lasso ne tranche pas sur la répartition~: c'est le solveur qui décide.

Pour les moindres carrés ordinaires, $X^\top X$ est singulière, et toute paire de somme correcte s'ajuste aussi bien. Une pseudo-inverse renvoie la solution de norme minimale, qui répartit à parts égales, tandis qu'une inversion directe échoue ou renvoie des valeurs énormes de signes opposés.

Pour la régression ridge, la pénalité $b_1^2 + b_2^2$, à somme donnée, est minimale lorsque les deux coefficients sont égaux. Chaque coefficient vaut $x^\top y/(2x^\top x + \lambda)$, et leur somme vaut $x^\top y/(x^\top x + \lambda/2)$. Ainsi, doubler une colonne réduit de moitié sa pénalité effective.

Pour le lasso\index{lasso}, $|b_1| + |b_2|$ est la même pour toute répartition de même signe, si bien que toutes ces répartitions ont la même fonction objectif. La descente par coordonnées partant de zéro tend à concentrer tout sur la première colonne et à laisser la seconde à zéro. Cela dépend de l'ordre et du point de départ.

\iqwhy{Avec des signaux presque identiques, le comportement est similaire~: le choix du lasso change d'un jour à l'autre et crée du turnover sans signification économique. La régression ridge, ou un regroupement des signaux, est plus stable.}
\end{solution}

\enlargethispage{-2\baselineskip}%
\begin{solution}{iq:iv1:regression:14}
Le lasso donne $(2, 0, 0.2)$, et la régression ridge donne $(1.5, -0.25, 0.6)$.

Avec des colonnes orthonormées, la fonction objectif se décompose en un problème par coefficient. À constante près, $\tfrac12\lVert y - Xb\rVert^2 = \tfrac12\sum_j (b_j - \hat\beta_j)^2$, où $\hat\beta_j$ est le coefficient des moindres carrés ordinaires.

En minimisant $\tfrac12(b - \hat\beta)^2 + \tfrac{\lambda}{2}b^2$, pour la régression ridge, on obtient $b = \hat\beta/(1 + \lambda)$~: chaque coefficient est divisé par 2.

Pour le lasso, on minimise $\tfrac12(b - \hat\beta)^2 + \lambda|b|$. Si $b > 0$, la dérivée s'annule en $b = \hat\beta - \lambda$, ce qui exige $\hat\beta > \lambda$. Le cas $b < 0$ est symétrique. Si $|\hat\beta| \le \lambda$, le point anguleux en zéro est le minimum. C'est le seuillage doux\index{seuillage doux},
\[
b = \mathrm{sign}(\hat\beta)\,\bigl(|\hat\beta| - \lambda\bigr)_+ ,
\]
représentée dans la \cref{fig:iv1:shrink}. Ainsi $3 \to 2$, $-0.5 \to 0$ et $1.2 \to 0.2$.

\begin{omfigure}
\begin{tikzpicture}
\begin{axis}[width=8cm,height=5.6cm,xmin=-3.3,xmax=3.3,ymin=-3.3,ymax=3.3,axis lines=middle,
  xlabel={\small MCO},ylabel={\small estimation},xlabel style={anchor=west},ylabel style={anchor=south},
  xtick={-3,-1,1,3},xticklabels={,,$1$,$3$},ytick={-2,2},tick label style={font=\footnotesize},clip=false]
\addplot[gray,dashed,domain=-3:3] {x};
\addplot[omThm,thick,samples=2,domain=-3:3] {x/2};
\addplot[omDef,very thick,samples=61,domain=-3:3] {x>1 ? x-1 : (x<-1 ? x+1 : 0)};
\node[font=\footnotesize,anchor=south] at (axis cs:-1,0) {$-1$};
\node[font=\footnotesize,anchor=south] at (axis cs:-3,0) {$-3$};
\node[font=\footnotesize,anchor=west,omDef] at (axis cs:3.05,2) {lasso};
\node[font=\footnotesize,anchor=west,omThm] at (axis cs:3.05,1.5) {ridge};
\node[font=\footnotesize,anchor=west,gray] at (axis cs:3.05,3) {MCO};
\end{axis}
\end{tikzpicture}

\omcaption{L'estimation en fonction du coefficient MCO, avec $\lambda = 1$ et des colonnes orthonormées. La régression ridge divise par 2 partout. Le lasso met à zéro tout ce qui se trouve dans $[-1, 1]$ et retranche 1 au reste, si bien que les grands coefficients conservent un biais fixe.}\label{fig:iv1:shrink}
\end{omfigure}

\iqfollow{Pourquoi n'existe-t-il pas de forme close en général~?}{La valeur absolue n'est pas dérivable, et des colonnes non orthogonales couplent les coefficients. On le résout par descente par coordonnées, qui applique cette même règle un coefficient à la fois. On peut aussi l'écrire comme un programme quadratique, en posant $b = b^+ - b^-$ avec $b^+, b^- \ge 0$ et en prenant $\sum(b^+ + b^-)$ comme pénalité.}
\iqfollow{Et sans le $\tfrac12$ devant l'erreur quadratique~?}{Le seuil devient $\lambda/2$.}
\end{solution}

\enlargethispage{-3\baselineskip}%
\begin{solution}{iq:iv1:regression:15}
Les coefficients sont sans biais sous linéarité, exogénéité et plein rang. Les erreurs types usuelles supposent aussi des erreurs de variance constante et non corrélées. Des erreurs normales ne sont pas nécessaires pour les grands échantillons.

Pour l'absence de biais, il faut $y = X\beta + \varepsilon$, $E[\varepsilon \mid X] = 0$ et un rang plein en colonnes. Pour des erreurs types correctes et la propriété de meilleur estimateur linéaire sans biais\index{théorème de Gauss--Markov}, il faut aussi des erreurs homoscédastiques et non corrélées. La normalité ne compte que pour les tests exacts sur petits échantillons.

Sur les rendements quotidiens d'actions~:
\begin{itemize}[leftmargin=*,itemsep=1pt,topsep=2pt]
\item La volatilité varie. Les coefficients restent sans biais, mais les erreurs types sont fausses. Utiliser des erreurs types robustes (de White), ou les normaliser par une prévision de volatilité.
\item Les erreurs sont corrélées, surtout avec des horizons qui se chevauchent. Utiliser les erreurs types de Newey--West ou de Hansen--Hodrick, ou des données sans chevauchement.
\item Queues épaisses. Quelques jours dominent l'ajustement. Le niveau du test tient raisonnablement avec des milliers de points, mais vérifier l'influence, et envisager une régression à perte de Huber ou la winsorisation.
\item L'exogénéité échoue à cause des facteurs omis et du look-ahead dans les variables. Celle-ci crée un biais, et davantage de données ne le corrigeront pas. Contrôler les facteurs et horodater strictement.
\item Les relations dérivent. Ajuster sur des fenêtres glissantes ou à pondération exponentielle, et tester la stabilité sur des sous-périodes.
\item Les régresseurs sont bruités. Les pentes tendent vers zéro.
\item Les régresseurs sont presque colinéaires. Les coefficients sont instables. Utiliser la régression ridge ou les combiner.
\end{itemize}

\iqwrong{«~Les erreurs doivent être normales pour que les MCO soient sans biais.~» La normalité n'a rien à voir avec le biais.}
\end{solution}

\begin{solution}{iq:iv1:regression:16}
Le $R^2$ attendu vaut $20/99 \approx 0.20$, le $R^2$ ajusté attendu vaut 0, et environ un coefficient possède $|t| > 2$.

Sous l'hypothèse nulle, $R^2$ suit une loi Bêta$(p/2,\ (n - p - 1)/2)$ de moyenne $p/(n - 1)$\index{$R^2$ ajusté}. Chacun des 20 régresseurs consomme un des 99 degrés de liberté au-delà de la moyenne, et chacun capte par hasard une part égale de la variance.

L'ajustement $1 - (1 - R^2)(n - 1)/(n - p - 1)$ a été construit pour être nul en moyenne dans ce cas.

Chaque statistique $t$ a 79 degrés de liberté, et $P(|t| > 2) \approx 0.049$, donc le nombre attendu est $20 \cdot 0.049 \approx 1$.

\iqwhy{Le $R^2$ dans l'échantillon (in-sample) de 20~\% obtenu sur du bruit indique ce qu'on peut attendre d'un ajustement avec de nombreux signaux et peu de données. Il faut comparer vos résultats à cette référence, et juger sur les résultats hors échantillon (out-of-sample).}
\end{solution}

\begin{solution}{iq:iv1:regression:17}
Avec 250 pas, environ 85~\% des régressions sont «~significatives~», avec un $R^2$ moyen d'environ $0.24$. Avec 1\,000 pas, la part monte à environ 93~\%, tandis que le $R^2$ reste au même niveau. Il faut utiliser des rendements, ou tester la cointégration.

Il s'agit d'une régression fallacieuse\index{régression fallacieuse}. Les deux séries dérivent, et deux séries qui dérivent semblent liées sur une fenêtre finie. La théorie standard suppose des erreurs stationnaires, et ici les résidus suivent eux-mêmes une marche aléatoire, si bien que la statistique $t$ n'a pas de loi limite. Elle croît comme $\sqrt{T}$~: passer de 250 à 1\,000 pas double sa dispersion. La pente ne converge pas, et elle reste aléatoire même avec des données infinies. Ces chiffres proviennent d'une simulation à graine fixée.

La différenciation y remédie. Régresser des rendements sur des rendements donne $|t| > 1.96$ environ 5~\% du temps, comme il se doit, avec un $R^2$ proche de zéro.

\iqfollow{Dans quels cas une régression en niveaux est-elle pertinente~?}{Lorsque les deux séries sont cointégrées\index{cointégration}, de sorte qu'une combinaison linéaire des deux est stationnaire. Les résidus ne dérivent alors pas, et c'est ce qu'on teste~: par exemple un test de racine unitaire sur les résidus, avec les valeurs critiques adéquates.}
\iqwhy{Deux graphiques de cours qui évoluent ensemble ne constituent pas une preuve. Une paire de trading a besoin d'un spread stationnaire, et les tests sur les niveaux de prix paraîtront excellents pour presque n'importe quelle paire.}
\end{solution}

\begin{solution}{iq:iv1:regression:18}
Ridge, sauf si l'on a une raison de croire que quelques signaux portent presque toute l'information.

Le lasso suppose la parcimonie, et la sélection de variables qu'il opère se paie par une instabilité~: parmi des signaux corrélés, il en retient un au hasard et change de choix d'un échantillon à l'autre, ce qui ajoute du turnover. La régression ridge\index{régression ridge} les conserve tous et réduit les coefficients dans les directions de faible variance. Avec $p/n = 0.4$, de nombreux petits effets et un signal faible, la réduction de la variance compte le plus.

Dans une simulation à graine fixée avec 500 signaux, corrélés à travers dix facteurs, et 1\,250 observations, où tous les signaux ont un petit effet aléatoire, le $R^2$ hors échantillon (out-of-sample) vaut environ 1.1~\% pour ridge et environ 0.8~\% pour lasso, chacun avec sa meilleure pénalité. Avec seulement 5 signaux forts et 495 inutiles, le lasso prend l'avantage.

Il faut revoir ce choix si l'on a de bonnes raisons de croire à la parcimonie, si l'on recherche l'interprétabilité ou moins d'entrées à maintenir, ou si chaque signal coûte de l'argent à faire tourner. L'elastic net est un compromis raisonnable.

\iqfollow{À quoi devez-vous faire attention~?}{Il faut standardiser les signaux, ne pas pénaliser l'ordonnée à l'origine, et choisir la pénalité par validation croisée ordonnée dans le temps avec un écart. Des plis aléatoires laissent fuir l'information par l'autocorrélation.}
\end{solution}

\begin{solution}{iq:iv1:regression:19}
La pente vaut $0.6$, soit la moitié de $1.2$. Régresser sur les rendements de marché courant et décalés de plusieurs jours, puis additionner les coefficients, retrouve $1.2$, et un horizon de rendement plus long donne à peu près le même résultat.

Le mouvement de marché d'une journée apparaît dans le cours de l'action le premier jour ultérieur où elle est cotée, et chaque jour elle se traite avec la probabilité $1/2$. Une fraction $1/2$ du mouvement apparaît donc le jour même, une fraction $1/4$ un jour plus tard, $1/8$ deux jours plus tard, et ainsi de suite. Régresser sur les rendements de marché à plusieurs décalages donne donc des pentes $0.6$, $0.3$, $0.15$, et ainsi de suite, dont la somme est la vraie valeur $1.2$. Ignorer les décalages laisse $0.6$.\index{cours obsolètes}

L'autre remède consiste à allonger l'horizon de rendement, puisqu'une part plus petite du mouvement échappe à la fenêtre. Avec des rendements hebdomadaires dans ce modèle, seuls les mouvements de fin de semaine peuvent y échapper, et en moyenne $(31/32)/5$ du mouvement d'une semaine y échappe. La pente vaut donc environ $1.2\,(1 - 31/160) \approx 0.97$.

\iqwhy{Les valeurs peu liquides et les instruments cotés à des heures de clôture différentes affichent tous un bêta faible et une corrélation faible sur des données journalières. Les clôtures asiatique et américaine, par exemple, sont séparées de plusieurs heures. Il ne faut pas prendre cela pour de la diversification~: à des horizons plus longs, la corrélation réapparaît. Le même effet fait baisser la corrélation à des fréquences très élevées.}
\end{solution}

\begin{solution}{iq:iv1:regression:20}
L'erreur type vraie vaut $\sqrt{3/n}$, et la formule usuelle donne $\sqrt{1/n}$. Elle sous-estime donc la vérité d'un facteur $\sqrt3 \approx 1.73$.

L'erreur sur la pente vaut $\sum x_i e_i / \sum x_i^2$, avec $e_i = |x_i|\varepsilon_i$. Pour un grand nombre d'observations, sa variance vaut
\[
\frac{E[x^2 e^2]}{n\,(E[x^2])^2} = \frac{E[x^4]}{n} = \frac{3}{n} ,
\]
car $E[x^2 e^2] = E[x^2 \cdot x^2]E[\varepsilon^2]$ et le moment d'ordre quatre d'une loi normale centrée réduite vaut 3. La formule usuelle est $s^2/\sum x_i^2$ avec $s^2 \to E[e^2] = E[x^2] = 1$, elle donne donc $1/n$. Elle utilise le bruit moyen, et ne voit pas que le bruit est le plus fort précisément là où $x$ pèse le plus dans la pente.

Pour corriger les erreurs types, il faut utiliser l'estimateur sandwich de White\index{erreurs types de White}~: $\sqrt{\sum x_i^2 \hat e_i^2}/\sum x_i^2$ avec $x$ centré. Il estime directement le $3/n$ et coïncide avec la dispersion vraie dans une simulation.

Pour améliorer l'estimation, il faut utiliser les moindres carrés pondérés\index{moindres carrés pondérés}, qui pondèrent chaque point par l'inverse de sa variance de bruit, ici $1/x_i^2$. Dans ce cas, l'estimateur devient la moyenne des $y_i/x_i$, de variance $1/n$, soit un tiers de la variance des MCO.

\iqfollow{Chacune de vos observations est la moyenne d'un nombre différent $m_i$ de points de données sous-jacents. Quelles pondérations utilisez-vous~?}{$m_i$, puisque la variance d'une moyenne vaut $\sigma^2/m_i$.}
\iqfollow{Vous disposez de deux estimations indépendantes d'un même nombre, avec des erreurs types de 2 et 1. Comment les combinez-vous~?}{Avec des pondérations inverses des variances, soit $0.2$ et $0.8$, et une erreur type de $\sqrt{0.8} \approx 0.89$.}
\end{solution}

\begin{solution}{iq:iv1:regression:21}
Les réponses sont (a) un facteur 1, donc aucune sous-estimation, (b) environ $\sqrt{20} \approx 4.47$, et (c) environ $3.33$.

La formule usuelle suppose que $x_t e_t$ est non corrélé dans le temps. En réalité, la variance de la pente est plus grande d'un facteur
\[
1 + 2 \sum_{k \ge 1} \rho_x(k)\,\rho_e(k) ,
\]
où $\rho_x$ et $\rho_e$ sont les autocorrélations du régresseur et de l'erreur\index{observations qui se chevauchent}. Deux fenêtres de 20 jours décalées de $k$ jours partagent $20 - k$ jours, donc $\rho_e(k) = 1 - k/20$ pour $k < 20$.

(a) Si le signal est indépendant d'un jour à l'autre, $\rho_x(k) = 0$, le facteur vaut 1. Le chevauchement de $y$ seul ne nuit pas à l'erreur type de la pente. Beaucoup de candidats passent à côté de ce point.

(b) Si le signal est presque constant, $\rho_x(k) \approx 1$, le facteur vaut $\sum_{|k|<20}(1 - |k|/20) = 20$, donc les erreurs types sont trop petites d'un facteur $\sqrt{20}$.

(c) Avec $\rho_x(k) = 0.9^k$, le facteur vaut $1 + 2\sum_{k=1}^{19} 0.9^k(1 - k/20) \approx 11.1$, dont la racine carrée vaut $3.33$. Les simulations reproduisent (a) et (c), et (b) est la limite de la même formule.

La solution consiste à utiliser les erreurs types de Newey--West\index{erreurs types de Newey--West} ou de Hansen--Hodrick, avec un nombre de retards au moins égal au chevauchement. Les pondérations de Newey--West décroissent avec le retard, il faut donc en prendre davantage~: dans la simulation avec 40 retards, on reste à 20~\% près de la vérité pour le cas (c).

\iqfollow{Qu'en est-il des erreurs simplement autocorrélées, sans chevauchement~?}{Avec un AR(1) dans le régresseur comme dans l'erreur, de coefficient $0.9$ chacun, le facteur vaut $(1 + 0.81)/(1 - 0.81) = 9.53$, donc les erreurs types sont trop petites d'un facteur $3.09$. La statistique de Durbin--Watson, qui vaut environ $2(1 - \rho)$, constitue le contrôle rapide sur les résidus.}
\end{solution}

\begin{solution}{iq:iv1:regression:22}
Le $R^2$ peut prendre toute valeur entre $0.36$ et 1. Et non, $x_1$ n'est pas inutile du tout~: avec les deux régresseurs, le $R^2$ vaut $0.842$, contre $0.16$ pour $x_2$ seul.

Pour deux régresseurs de corrélations $r_1$, $r_2$ avec $y$ et $\rho$ entre eux,
\[
R^2 = \frac{r_1^2 + r_2^2 - 2 r_1 r_2 \rho}{1 - \rho^2} .
\]
Avec $r_1 = r_2 = 0.6$, cela vaut $0.72/(1 + \rho)$. Quand $\rho \to 1$, les deux régresseurs sont la même variable et $R^2 \to 0.36$. Abaisser $\rho$ fait monter $R^2$. La matrice de corrélation de $y, x_1, x_2$ doit être semi-définie positive, ce qui impose ici $\rho \ge 2 \cdot 0.36 - 1 = -0.28$. Pour cette valeur, $R^2 = 1$.

Dans le second cas, avec $r_1 = 0$, $r_2 = 0.4$ et $\rho = 0.9$, $R^2 = 0.16/0.19 \approx 0.842$. C'est une variable suppressive\index{variable suppressive}~: $x_1$ n'a aucun lien direct avec $y$, mais il annule la part de $x_2$ qui est du bruit pour la prédiction de $y$.

\iqwrong{«~Le $R^2$ conjoint est au plus égal à la somme des $R^2$ séparés~», ou «~une variable de corrélation nulle avec $y$ ne peut pas aider~». Les deux affirmations sont fausses.}
\iqwhy{Il ne faut pas sélectionner les signaux un à un selon leur corrélation avec la cible. Une couverture repose sur la même idée~: un instrument non corrélé à la cible peut tout de même réduire le risque lorsqu'on détient une autre position.}
\end{solution}

\begin{solution}{iq:iv1:regression:23}
La meilleure combinaison atteint $0.0577$, $0.0707$ et $0.1$ dans les trois cas. Avec des corrélations de $0.05$ et $0.02$ et $\rho = 0.5$, la pondération du second signal est négative, alors qu'il prédit positivement pris isolément.

La corrélation de la meilleure combinaison linéaire avec la cible vaut $\sqrt{R^2}$\index{combinaison de signaux}, et la formule pour deux régresseurs donne $R^2 = 2r^2/(1 + \rho)$ avec des corrélations égales. Pour $r = 0.05$~:
\[
\rho = 0.5 : 0.0577, \qquad \rho = 0 : 0.0707, \qquad \rho = -0.5 : 0.1.
\]
La diversification joue~: des signaux corrélés apportent peu, des signaux indépendants s'additionnent en $\sqrt2$, et des signaux négativement corrélés, de même signe de prédiction, s'additionnent davantage.

Les pondérations optimales valent $w \propto \Sigma^{-1}r$. Avec $\Sigma = \bigl(\begin{smallmatrix}1 & 0.5\\ 0.5 & 1\end{smallmatrix}\bigr)$ et $r = (0.05, 0.02)$,
\[
w = \frac{1}{0.75}\bigl(0.05 - 0.5 \cdot 0.02,\ \ 0.02 - 0.5 \cdot 0.05\bigr) = (0.0533,\ -0.0067).
\]
La corrélation propre du second signal, de $0.02$, est inférieure aux $0.025$ qu'il obtiendrait simplement par sa corrélation avec le premier. Il est donc vendu à découvert par rapport au premier, ce qui ajoute un peu~: la corrélation vaut $0.0503$ contre $0.05$ seul.

\iqwhy{C'est la régression du rendement sur les signaux, et elle montre pourquoi combiner vaut mieux que sélectionner. Le gain de $0.05$ à $0.0577$ est modeste mais réel~: un ratio de Sharpe annuel d'environ $0.92$, contre $0.79$. Il suffit de peu pour que le signe d'une pondération s'inverse, d'où l'intérêt de rétrécir les pondérations.}
\end{solution}

\begin{solution}{iq:iv1:regression:24}
La meilleure valeur est $\lambda = \sigma^2/\beta^2 = 4$, qui donne une erreur quadratique moyenne de $0.2$ contre 1 pour les MCO, soit une réduction de 80~\%.

L'estimation ridge $\hat\beta/(1 + \lambda)$ a pour biais $-\lambda\beta/(1 + \lambda)$ et pour variance $\sigma^2/(1 + \lambda)^2$, donc
\[
\mathrm{MSE}(\lambda) = \frac{\sigma^2 + \lambda^2\beta^2}{(1 + \lambda)^2}.
\]
Annuler la dérivée donne $\lambda^* = \sigma^2/\beta^2$, et le minimum vaut $\sigma^2\beta^2/(\sigma^2 + \beta^2)$. Avec $\sigma = 1$ et $\beta = 0.5$, on obtient $\lambda^* = 4$ et $\mathrm{MSE} = 0.2$. L'estimation vaut $\hat\beta/5$\index{régression ridge}~: on jette les quatre cinquièmes.

Plus le signal est faible par rapport au bruit, plus il faut rétrécir. Un signal de $t$-statistique vraie de $0.5$ doit être fortement réduit, alors qu'un signal avec $\beta = 5$ n'est réduit que de $\lambda^* = 0.04$. Bien sûr, $\beta$ est inconnu, et c'est pourquoi on choisit $\lambda$ par validation. Mais la conclusion tient~: dans les situations de faible rapport signal sur bruit, les estimations sans biais sont de mauvais paris.

\iqfollow{Quelle est la lecture bayésienne~?}{Avec un a priori $\beta \sim N(0, \tau^2)$, la moyenne a posteriori est l'estimation ridge avec $\lambda = \sigma^2/\tau^2$. Pour $\tau^2 = 0.25$, cela multiplie à nouveau l'estimation par $1/5$.}
\iqfollow{Et avec beaucoup de régresseurs~?}{Dans la base des valeurs singulières, chaque direction est multipliée par $d^2/(d^2 + \lambda)$, si bien que les directions où $X$ a peu de variance sont les plus rétrécies.}
\end{solution}

\begin{solution}{iq:iv1:regression:25}
Garder $x > 0$ laisse la pente à $0.5$ et ramène la corrélation à environ $0.33$. Garder $y > 0$ ramène la pente à environ $0.22$.

Sélectionner sur $x$ laisse la relation intacte\index{biais de sélection}~: à $x$ donné, $y$ a toujours pour moyenne $0.5\,x$, donc la pente reste $0.5$. Seule la dispersion de $x$ diminue. La variance d'une normale centrée réduite au-dessus de zéro vaut $v = 1 - 2/\pi \approx 0.363$. La corrélation baisse parce que la variance expliquée baisse~:
\[
\text{corr}^2 = \frac{\rho^2 v}{\rho^2 v + 1 - \rho^2} = \frac{0.25 \cdot 0.363}{0.25 \cdot 0.363 + 0.75} \approx 0.108, \quad \text{corr} \approx 0.33.
\]

Sélectionner sur $y$ casse la relation, car le bruit compte désormais~: pour un $x$ faible, on ne garde que les points dont l'erreur est grande et positive. Les erreurs ne sont donc plus indépendantes de $x$, et la pente baisse. En utilisant la relation gaussienne, avec $x = \rho y + \sqrt{1 - \rho^2}\,z$ et la même variance $v$ pour $y$ au-dessus de zéro, on obtient $\mathrm{Cov}(x,y) = \rho v$ et $\mathrm{Var}(x) = \rho^2 v + 1 - \rho^2$, donc
\[
\text{pente} = \frac{\rho v}{\rho^2 v + 1 - \rho^2} = \frac{0.5 \cdot 0.363}{0.841} \approx 0.216.
\]

\iqwhy{Filtrer sur une variable que l'on utilise comme régresseur est sans danger pour la pente. Filtrer sur le résultat ne l'est pas, et cela vaut aussi pour le fait de ne garder que les survivants, que les ordres exécutés ou que les transactions rentables. Par exemple, une étude portant sur les «~actions qui ont monté~» ou les «~ordres exécutés~» sera biaisée, et aucune quantité de données n'y changera rien.}
\end{solution}

\begin{solution}{iq:iv1:regression:26}
Son levier vaut $0.101$, le retirer abaisse la pente d'environ $0.022$, et son erreur de prévision hors échantillon est d'environ $2.22$.

Le levier\index{levier} est la diagonale de la matrice chapeau, c'est-à-dire de combien l'observation déplace sa propre valeur ajustée. Avec une constante,
\[
h = \frac1n + \frac{(x - \bar x)^2}{\sum (x_i - \bar x)^2} = \frac{1}{1000} + \frac{100}{1000} = 0.101.
\]
Les leviers somment au nombre de coefficients, donc la moyenne ici vaut $2/1000$, et ce point en a plus de 50 fois autant.

Supprimer un point admet des formules exactes\index{formule de validation croisée leave-one-out}, donc aucun réajustement n'est nécessaire~:
\[
\text{erreur de prévision sans lui} = \frac{e}{1 - h} = \frac{2}{0.899} \approx 2.22,
\]
\[
\text{variation de la pente} = \frac{(x - \bar x)\,e}{(1 - h)\sum(x_i - \bar x)^2} = \frac{10 \cdot 2}{0.899 \cdot 1000} \approx 0.022.
\]
\iqfollow{Pourquoi cela compte-t-il du point de vue du calcul~?}{La même formule donne l'erreur complète de validation croisée leave-one-out à partir d'un seul ajustement, en divisant chaque résidu par $1 - h_i$.}
\iqwhy{Une observation de grand $x$ peut à elle seule piloter la pente, et son résidu est petit parce que la droite s'est déplacée vers elle. Il faut donc regarder le levier et pas seulement les résidus. Sur des données de rendements, cela signifie les jours de krach. Il faut vérifier que le signe tient sans eux, ou borner le régresseur lui-même, par winsorisation de $x$ ou au moyen d'un ajustement à influence bornée. Une perte de Huber seule ne suffit pas, puisqu'elle plafonne les grands résidus, et un point de fort levier peut en avoir un petit.}
\end{solution}

\begin{solution}{iq:iv1:regression:27}
Les MCO donnent $4/3$, et non 1. L'instrument donne $\mathrm{Cov}(z, y)/\mathrm{Cov}(z, x) = 1$. Un instrument faible donne une estimation très large, biaisée vers les MCO.

Les MCO ont pour pente $\mathrm{Cov}(x, y)/\mathrm{Var}(x)$. Ici $\mathrm{Var}(x) = 3$, et $\mathrm{Cov}(x, y) = \mathrm{Var}(x) + \mathrm{Cov}(x, u) = 3 + 1 = 4$, soit $4/3$~: $x$ et l'erreur évoluent ensemble, et la régression attribue l'effet à $x$.

La variable $z$ est un instrument\index{variables instrumentales}~: elle déplace $x$, et elle n'a aucune autre voie vers $y$. Alors $\mathrm{Cov}(z, y) = \beta\,\mathrm{Cov}(z, x)$, puisque $z$ est non corrélée avec $u$. Donc $\beta = \mathrm{Cov}(z, y)/\mathrm{Cov}(z, x) = 1/1 = 1$. Les moindres carrés en deux étapes donnent le même nombre~: on régresse $x$ sur $z$, puis $y$ sur les valeurs ajustées. Mais si on le fait à la main, les erreurs types de la seconde étape sont fausses, car elles utilisent les mauvais résidus.

Avec un instrument faible, le dénominateur est proche de zéro et bruité, et une petite erreur sur celui-ci déplace beaucoup le rapport. Dans une simulation avec $x = 0.1\,z + u + v$ et 200 observations, l'estimation médiane est tirée nettement au-dessus de 1, vers la valeur des MCO, et la dispersion est très large. Il faut toujours regarder la régression du premier étage.

\iqwhy{Une régression des rendements sur le flux d'ordres souffre de ce problème~: les prix et le flux réagissent à la même nouvelle. Estimer l'impact sur les prix par des MCO simples le surestime, puisque le flux informé s'accompagne de mouvements qu'il n'a pas causés. Il faut utiliser comme instrument un flux indépendant des nouvelles, comme les transactions liées aux rééquilibrages d'indices ou un calendrier d'exécution randomisé.}
\end{solution}

\begin{solution}{iq:iv1:regression:28}
On s'attend à environ $-1~\%$, un nombre négatif. Il faut environ 1\,980 observations avant que le $R^2$ hors échantillon n'atteigne zéro.

Avec $p$ coefficients estimés et $n$ observations, chaque coefficient comporte une erreur d'estimation, et l'erreur de prévision croît en conséquence\index{$R^2$ hors échantillon}. Avec des régresseurs bien conditionnés, l'erreur quadratique moyenne hors échantillon vaut approximativement la variance du bruit multipliée par $(1 + p/n)$. Avec une variance totale de 1, la variance du bruit vaut $1 - R^2$, donc
\[
R^2_{\text{oos}} \approx R^2 - \frac{p}{n}(1 - R^2) = 0.01 - 0.02 \cdot 0.99 \approx -0.0098 .
\]
Le $R^2$ en échantillon du même ajustement vaut environ 3~\%, puisqu'il ajoute $p/n$ au lieu de le soustraire. Annuler la valeur hors échantillon donne $n \approx p(1 - R^2)/R^2 = 1980$. Pour un seul régresseur, le seuil de rentabilité n'est que de 99 observations.

\iqwrong{«~La vraie relation est positive, donc l'ajustement doit battre la moyenne.~» Un signal réel mais faible, estimé avec trop de coefficients, perd face à la moyenne simple. Son erreur d'estimation est plus grande que le signal.}
\iqwhy{C'est la raison la plus fréquente d'un backtest qui échoue en trading réel alors qu'aucun bug n'est en vue. Le remède consiste à réduire $p$ ou à rétrécir~: la régression ridge coûte un peu de biais et supprime une grande part du terme $p/n$.}
\end{solution}

\enlargethispage{-2\baselineskip}%
\begin{solution}{iq:iv1:regression:29}
La pente moyenne vaut environ $+0.03$, alors que la vraie valeur est 0. Le test rejette beaucoup trop souvent~: environ 16~\% du temps, et non 5~\%.

C'est le biais de Stambaugh\index{biais de Stambaugh}. Le régresseur $x_t$ est non corrélé avec le choc suivant, mais il n'est pas exogène~: il dépend des valeurs passées de $u$ par l'intermédiaire de $v$. Écrivons $u_t = \gamma v_t + w_t$, avec $w$ indépendant de $v$ et $\gamma = \mathrm{Cov}(u, v)/\mathrm{Var}(v) = -0.9$. La pente de $v_{t+1}$ sur $x_t$ vaut exactement $\hat\rho - \rho$, donc
\[
\hat\beta - \beta = \gamma(\hat\rho - \rho) + (\text{pente sur } w),
\]
et le dernier terme est de moyenne nulle.

L'erreur sur $\hat\beta$ est donc l'erreur sur l'estimation de la persistance, multipliée par $\gamma$. La persistance est sous-estimée sur les petits échantillons\index{biais de Kendall}~: l'approximation standard donne $E[\hat\rho - \rho] \approx -(1 + 3\rho)/T = -3.85/120 \approx -0.032$. D'où $E[\hat\beta] \approx (-0.9)(-0.032) \approx 0.029$. Une simulation à graine fixée donne une pente moyenne de $0.032$, un biais de $\hat\rho$ de $-0.035$ et une $t$-statistique moyenne d'environ $0.7$. Le test rejette environ 16~\% du temps. On peut le chiffrer au tableau. Avec $\mathrm{Var}(x) = 1/(1 - \rho^2) \approx 10.3$, l'erreur type vaut au moins environ $1/(3.2\sqrt{120}) \approx 0.029$, soit autant que le biais, si bien que la $t$-statistique se décale vers le haut d'au plus environ 1. Dans une approximation gaussienne, un décalage de $0.7$ à 1 fait passer le taux de rejet de 5~\% à une valeur comprise entre $17\%$ et $26\%$. Les 16~\% simulés sont légèrement en dessous de cette fourchette, car le décalage moyen est à l'extrémité basse et la dispersion de la statistique est un peu inférieure à 1.

\iqwhy{C'est le cadre type de la prévision des rendements à l'aide d'un ratio de valorisation, tel que le rendement du dividende~: une hausse du prix fait monter le rendement et baisser le ratio, si bien que $\gamma$ est fortement négatif, et que le ratio est très persistant. Une part des preuves apportées par une telle régression relève de ce biais. On le réduit en corrigeant par une estimation de $\gamma(\hat\rho - \rho)$, ou en testant avec une distribution nulle obtenue par simulation.}
\iqfollow{Le même biais apparaît lorsqu'on estime un retour à la moyenne. Que produit-il~?}{Il surestime la vitesse du retour. Avec $\rho = 0.95$, la vraie demi-vie vaut $13.5$ périodes. L'estimation moyenne sur 120 points, $\hat\rho \approx 0.915$, donne $7.8$.}
\end{solution}

\begin{solution}{iq:iv1:regression:30}
Les coefficients valent environ $0.42$ sur $w$ et environ $0.47$ sur $x_2$. On conclurait que $x_2$ a un effet fort, et c'est faux. Son coefficient ne reflète que le bruit de $w$.

Résolvons les équations normales dans la population\index{erreurs sur les variables}. Posons $\rho = 0.8$ et $s^2 = 0.5$. Alors $\mathrm{Var}(w) = 1 + s^2$, $\mathrm{Cov}(w, x_2) = \rho$, $\mathrm{Cov}(w, y) = 1$ et $\mathrm{Cov}(x_2, y) = \rho$, ce qui donne
\[
\begin{pmatrix} 1 + s^2 & \rho \\ \rho & 1 \end{pmatrix}\begin{pmatrix} b_1 \\ b_2 \end{pmatrix} = \begin{pmatrix} 1 \\ \rho \end{pmatrix}.
\]
La solution est
\[
b_1 = \frac{1 - \rho^2}{1 + s^2 - \rho^2} = \frac{0.36}{0.86} \approx 0.42, \qquad b_2 = \frac{\rho\,s^2}{1 + s^2 - \rho^2} = \frac{0.4}{0.86} \approx 0.47.
\]
Lorsque $s^2 = 0$, le coefficient sur $x_2$ est nul, comme il se doit. Quand $s^2$ croît, la régression fait moins confiance à $w$ et s'appuie davantage sur $x_2$, car $x_2$ est un substitut propre de $x_1$. À la limite, $b_2$ tend vers $\rho = 0.8$, le coefficient de $y$ sur $x_2$ seul.

Davantage de données n'aident pas. Le biais subsiste pendant que l'erreur type diminue, si bien que la $t$-statistique croît sans limite. Avec 2\,500 observations, elle dépasse nettement 8 dans une simulation.

\iqwrong{Croire que le bruit sur un régresseur n'affecte que le coefficient de ce régresseur. Le biais se propage à tout régresseur corrélé.}
\iqwhy{C'est ainsi qu'un nouveau facteur peut «~survivre aux contrôles~» sans rien apporter. Si le facteur existant que l'on contrôle est mesuré avec bruit, comme une valeur comptable périmée ou un bêta bruité, toute variable qui lui est corrélée paraîtra significative.}
\end{solution}
